%% Section V-A and V-B draft: protocol and baselines.
%% The dataset prose carries over from the previous manuscript nearly verbatim
%% (it survived four red-team rounds); the protocol paragraphs are rewritten for
%% the incremental estimator and the causal/committed reporting rule.
%% MPF numbers: research/mpf_*_results.csv, best per case, no-NN columns only.
%% NOTE: the sigma_beta=100 MPF rerun (research/mpf_prior.jl) may upgrade the
%% MPF paragraph from "unstable under best tuning" to a counted baseline; the
%% paragraph is written so only its second half changes in that case.

\section{Experiments and Results}\label{sec:results}

\subsection{Dataset, maps, and protocol}\label{sec:protocol}
All experiments use the public SGL~2020 challenge dataset~\cite{gnadt2020}: a
Cessna~208B instrumented with a flight-grade INS, several onboard scalar
magnetometers (a compensated tail-stinger sensor and uncompensated cabin sensors
Mag~4 and Mag~5, all optically pumped cesium-vapor units), a three-axis
fluxgate, and reference truth. The flights serve distinct purposes, and we
choose lines accordingly (Table~\ref{tab:lines}). Flt1007 is the free-flight
mission, in which the aircraft flies freely within the mapped areas, so its
lines are the closest to a navigation task and are what prior
work~\cite{gnadt2020,canciani2017,hager2026} reports; 1007.06 is our
representative line for that reason. Flt1003 is a crustal-field survey flight,
and its lines 1003.02 and 1003.08 are straighter than the free-flight legs.
Flt1006 is a calibration-manoeuvre flight; its single short line is listed for
completeness but not counted, since at \SI{14}{\minute} the warm-up leaves too
little line to score. The map-building flights (Flt1004--1005) are not used.
Two honest points remain. None of these lines is a dedicated long,
straight-and-level operational transit; the free-flight and survey legs carry
attitude variation, which stresses the compensation but also aids its
observability by exciting the regressor, so performance on a pure operational
cruise is not established here. And the lines span two altitude regimes, four
near the \SI{400}{\meter} survey-drape altitude and 1007.02 near
\SI{800}{\meter}, where upward continuation attenuates the anomaly.

Map matching uses the Eastern and Renfrew high-resolution anomaly grids,
upward-continued to flight altitude, with the IGRF core field restored so that
measurement and map share an absolute scale. Cold start means the compensation
coefficients are initialized at zero with no calibration flight; the navigation
solution itself is accurately initialized (position standard deviation
\SI{0.1}{\meter}, Table~\ref{tab:params}), so the term denotes a compensation
cold start, not an unknown-position start.

The error metric is the horizontal distance root-mean-square (DRMS) of position
error, counted after a \SI{600}{\second} warm-up. Following
Section~\ref{sec:incremental}, the proposed estimator is reported twice in every
table: the \emph{causal} value, read at the newest state with no look-ahead,
and the \emph{committed} value, read as states leave the \SI{300}{\second} lag.
The causal value is the like-for-like comparison against the recursive
baselines, which also use no future data; the committed value is the output of
a system permitted to lag real time by the lag length, and the look-ahead
behind every reported number is stated in its table. One configuration
(Table~\ref{tab:params}) is used for every line and both magnetometers; nothing
is tuned per case, and parameter sensitivity is quantified in the ablation of
Section~\ref{sec:ablation}. All numbers are produced by scripted runs under a
fixed software environment and are available from the authors on request.

\subsection{Baselines}\label{sec:baselines}
Four references frame the comparison, all run by us on the same lines, the same
maps, and the same inertial model.

\emph{Free-inertial INS.} The unaided INS drift (Table~\ref{tab:lines}) is the
floor every aided estimator must beat, and the common reference that makes DRMS
values comparable across lines of different length and roughness.

\emph{Online-TL EKF.} The direct recursive counterpart of the proposed
estimator: the same Pinson error state augmented with the same \num{19}
time-varying TL coefficients, updated by the same scalar measurement, in an
extended Kalman filter~\cite{gnadt2022}. It shares every model term with the
factor graph, so the margin between the two isolates the estimator structure,
recursive commitment against windowed re-linearization, from the modelling.

\emph{Online EKF+TL+NN.} The strongest published causal method for this
dataset augments the TL state with the weights of a small feed-forward network
that absorbs what the linear basis cannot~\cite{gnadt2022,hager2026}. We run
the reference implementation of~\cite{gnadt2022}, tuned by us for a stable cold
start, and do not claim to reproduce the extended filter of~\cite{hager2026}.
Cold start forces one substantive change: with no calibration flight to
warm-start from, the network output normalization is initialized from the first
five minutes of onboard interference; a naive port instead drives the network
bias to $\sim\!5\times10^{4}$\,\si{\nano\tesla} and the filter to \texttt{NaN}.
We keep the uncompensated scalar out of the network inputs: a
position-dependent input carries the map gradient, so the network absorbs map
signal and position drifts while the residual stays small (the scalar-augmented
variant of~\cite{hager2026} escapes this only by freezing the network after its
transient). The tuned filter uses the permanent TL group with an eight-unit
hidden layer and reaches \num{48.8}/\SI{18.3}{\meter} on 1007.06 Mag~4/5,
consistent with published cold-start operating
points~\cite{gnadt2022,hager2026}.

\emph{Marginalized particle filter (MPF+TL).} The nonparametric reference,
included because particle methods are the standard resort when the map term is
strongly nonlinear~\cite{nordlund2009,park2024}. Since the TL model is linear
in its coefficients, $\boldsymbol{\beta}$ is conditionally linear-Gaussian and
fits the Rao--Blackwellized structure exactly: our implementation carries
position in \num{1000} particles and marginalizes the remaining states,
including $\boldsymbol{\beta}$, in the conditional Kalman block, with the
fluxgate row entering the linear measurement Jacobian. Weights are computed in
the log domain and position diversity is reinjected after resampling. Despite
this, and under per-case tuning of measurement variance, resampling threshold
and roughening that no other estimator in this paper receives, the filter is
unstable at a cold start: its best per-case results range from
\SI{38.7}{\meter} on the cleanest Mag~5 leg to kilometres on Mag~4, with strong
seed dependence. We report it in the text rather than in the comparison tables,
and we do not claim it represents the best achievable particle filter.
Attaching the neural network of the EKF+TL+NN to the particle filter does not
rescue it and is sharply harmful on the well-conditioned sensor, degrading
Mag~5 by one to two orders of magnitude across lines, the same direction seen
for the NN-augmented EKF where the linear basis already suffices.
